\section{Review Method and Corpus}
\label{sec:method}

\subsection{Review design}

The study is a curated scoping review. Scoping review methodology is appropriate
when concepts, evidence types, and evaluation practices are heterogeneous and the
goal is to map a field rather than estimate one pooled treatment
effect~\cite{tricco2018prismascr}. The accompanying \texttt{review\_protocol.md}
records the scope, eligibility criteria, evidence fields, update procedure, and
limitations. The machine-readable \texttt{evidence\_matrix.csv} is generated from
the workspace manifest by a versioned script; rule-assisted annotations remain
subject to manual audit. A second evidence layer, stored in
\texttt{precision\_annotations/}, contains full-text precision coding and an
individual evidence card for each of 30 purposively selected core papers.

The corpus was frozen on 12 August 2026. It began from a decision-oriented archive
assembled for research planning and was expanded around five intersecting areas:
(i) tactile and visuo-tactile representation learning; (ii) bimanual dexterous
manipulation and handover; (iii) graph, contact-field, object-centric, and
physics-informed models; (iv) tool--workpiece contact control; and (v) datasets
and benchmarks. Backward and forward chaining was used selectively for landmark
and directly competing works. Author project pages, arXiv records, official
conference programmes, proceedings pages, and publisher records were used to
disambiguate status.

The current draft follows the reporting intent of PRISMA-ScR but does not claim
formal PRISMA compliance. It does not reproduce database-by-database search
strings, use dual independent screening, or apply a validated cross-design
risk-of-bias tool. These omissions matter: the numerical corpus summary is an
audit of this archive, not a bibliometric description of the whole field.

\subsection{Two-layer evidence design}

The archive-wide layer covers all 74 records. Title, abstract, publication and
artifact metadata, and selectively consulted full text support eligibility,
status accounting, and coarse six-axis coding. This layer provides breadth but
is not used as if every row were an independently extracted full-text study.

The precision layer was completed on 13 August 2026 for 30 papers selected for
their decision relevance rather than for statistical representativeness. The
subset comprises 19 Level A records, coded through full text plus available
artifact or code evidence, and 11 Level B records, coded through full-text method
and evidence analysis. It covers direct system competitors, executable
benchmarks, graph and relational models, contact fields and anchors, mechanics
and control references, tool-use systems, and candidate local tactile encoders.
Each record distinguishes (i) the authors' claim, (ii) evidence located in the
paper or artifact, and (iii) reviewer inference. It also separates deployment
inputs from training-only privilege and codes object, sensor, embodiment, and
contact-topology OOD as different claims. All 30 cards contain evidence locations;
the procedure is full-text evidence coding, not independent reproduction.

\begin{table}[t]
\centering
\caption{Review questions and the evidence required to answer them.}
\label{tab:rqs}
\begin{tabularx}{\linewidth}{P{0.08\linewidth} Y Y}
\toprule
\textbf{RQ} & \textbf{Question} & \textbf{Primary evidence fields} \\
\midrule
RQ1 & What must be preserved for bimanual contact-rich manipulation? & Task topology, predicted quantities, temporal horizon, interface identity. \\
RQ2 & How is touch organized across sensors, bodies, entities, and time? & Representation scale, tokenization, spatial frame, structural prior. \\
RQ3 & Is the structure available at deployment? & Sensor inputs, privileged state, inferred versus supplied topology, future reference. \\
RQ4 & How is usefulness demonstrated? & Baselines, split unit, probes, control metrics, OOD tests, interventions, uncertainty. \\
RQ5 & Where are the actionable gaps? & Missing topology--method cells, weak evidence links, unavailable labels or benchmarks. \\
\bottomrule
\end{tabularx}
\end{table}

\subsection{Eligibility and coding}

Scholarly work was included when it materially addressed robot touch,
contact-state or dynamics estimation, tactile-informed control, bimanual
coordination, handover, tool use, or an enabling benchmark. Generic graph and
physics-learning papers were included only where their formulation directly
informs relational contact modelling. A community tactile-data index was retained
as a resource record but excluded from study-level synthesis. Duplicate versions
were represented by the most informative accessible record, with publication
status documented rather than silently upgraded.

Each record was coded along six axes:

\begin{enumerate}
  \item \textbf{Task/contact topology:} single interface, multi-finger hand,
  two hands sharing an object, sequential/overlapping handover, hand--tool--
  environment, or hand--tool--workpiece--hand.
  \item \textbf{Representation scale:} local patch, finger/link, hand,
  object/interface, cross-entity relation, or trajectory-level latent.
  \item \textbf{Structural prior:} token order, kinematics, object-centric
  coordinates, spatial anchors, contact fields/modes, or graph connectivity.
  \item \textbf{Supervision source:} instance discrimination, reconstruction,
  cross-modal alignment, action-conditioned prediction, simulator labels,
  force/pose measurement, or task reward.
  \item \textbf{Deployment observability:} directly sensed, perception-derived,
  training-only privilege, or unavailable future reference.
  \item \textbf{Evidence:} representation probe, open-loop prediction,
  closed-loop control, real transfer, OOD topology, intervention, and uncertainty.
\end{enumerate}

The axes are kept separate because they answer different validity questions.
A real-robot paper can use weak splits; a simulation paper can isolate a causal
mechanism cleanly; a high-capacity pretrained encoder can be deployable but offer
no cross-interface structure; and a physically interpretable graph can depend on
unavailable state.

\subsection{Corpus accounting and status discipline}

The archive contains 74 records: 12 core or directly competing systems, 19
bimanual and handover works, 16 tactile-representation works, 13 graph, physics,
and contact works, eight tool--workpiece works, and six dataset or benchmark records. Sixty-nine
PDFs are now locally validated; five records retain authoritative
source links without a local PDF. Seventy-three scholarly studies enter the
narrative synthesis and one community index is resource-only.

After correcting several manifest entries for identifiable formal proceedings,
the evidence matrix records 46 peer-reviewed works, eight works listed as accepted
or on an official programme, 19 preprints or author-claim records, and one
community resource. ``Accepted'' is not treated as equivalent to a final
publisher record. In particular, rapidly emerging 2026 works are discussed as
directions and evidence, not as settled consensus.

\begin{figure}[t]
\centering
\begin{tikzpicture}[
  box/.style={draw=reviewblue, rounded corners, fill=reviewlight, minimum width=31mm,
              minimum height=11mm, align=center, font=\small},
  status/.style={draw=reviewteal, rounded corners, fill=white, minimum width=31mm,
                 minimum height=10mm, align=center, font=\small},
  arr/.style={-{Latex[length=2mm]}, thick, draw=reviewgray}
]
\node[box] (archive) at (0,3.1) {74 archived\\records};
\node[box] (screen) at (4.5,3.1) {title/abstract and\\status audit};
\node[box] (synth) at (9.0,3.1) {73 scholarly works\\in synthesis};
\node[box, minimum width=39mm] (precision) at (6.65,1.35)
  {30 core papers\\full-text precision-coded};
\node[box] (resource) at (10.5,1.35) {1 community index\\resource-only};
\draw[arr] (archive) -- (screen);
\draw[arr] (screen) -- (synth);
\coordinate (fork) at (9.0,2.25);
\draw[arr] (synth.south) -- (fork) -| (precision.north);
\draw[arr] (fork) -| (resource.north);

\node[font=\scriptsize\bfseries, text=reviewgray] at (4.5,0.05)
  {Publication-status distribution of all 74 archived records};
\draw[reviewgray!45, thin] (-1.55,-0.22) -- (10.55,-0.22);
\node[status] (peer) at (-0.9,-1.0) {46\\peer-reviewed};
\node[status] (accept) at (2.7,-1.0) {8 accepted /\\programme-listed};
\node[status] (pre) at (6.3,-1.0) {19 preprints /\\author claims};
\node[status] (res) at (9.9,-1.0) {1\\resource};
\end{tikzpicture}
\caption{Two-layer corpus flow and publication-status accounting. All 74 records
receive archive-wide screening and coarse coding; the purposive 30-paper subset
receives full-text precision coding. Counts describe the curated archive, not the
global literature.}
\label{fig:corpus}
\end{figure}

\subsection{Synthesis strategy}

The synthesis is concept-centric rather than paper-by-paper. Representative
methods are compared within the same information problem, and negative space is
made explicit: for example, a work may demonstrate cross-sensor transfer without
testing cross-hand load inference, or closed-loop bimanual success without an
explicit tactile representation ablation. No pooled accuracy or success rate is
reported because task definitions, sensor density, action spaces, data scale,
embodiments, and test splits are not commensurate. Descriptive counts from the
precision subset are reported only when the denominator and coding rule are
explicit. In particular, zero of the 30 core papers evaluated strict held-out
contact-topology OOD under deployment-observable inference. Here ``strict'' means
that the participant/edge structure itself is held out; unseen objects, contact
modes within a fixed structure, tool geometry, or semantic task compositions do
not qualify. Claims that a specific intersection is comparatively sparse are
stated as inferences from the mapped corpus rather than field-wide publication
counts.
